\documentclass[]{article}


\usepackage{color,soul}   				% For highlighting with \hl{}

\newcommand{\nue}{\nu_{e}}
\newcommand{\numu}{\nu_{\mu}}
\newcommand{\nueb}{\bar{\nu}_{e}}
\newcommand{\numub}{\bar{\nu}_{\mu}}
\newcommand{\nub}{\bar{\nu}}


\begin{document}

\section{Analysis Strategy}

In the $\nueb$ appearance analysis, the significance of $\nueb$ appearance is given as a p-value by introducing the parameter $\beta$, which multiplies the $\numub \rightarrow \nueb$ oscillation probability by either 0 or 1. The data are compared to two null hypotheses: a) No $\nueb$ appearance ($\beta = 0$) and b) $\nueb$ appearance according to the current best knowledge of the PMNS matrix ($\beta = 1$), where both hypotheses fully account for uncertainties in the values of the oscillation and systematic parameters. Both rate-only and rate+shape analyses are performed on each hypothesis to obtain two different p-values for each hypothesis.

The total number of observed events in the $\nub$ mode single-ring e-like sample, without using any information from the reconstructed energy shape, is used to calculate the rate-only p-value. $\Delta \chi^2 = \chi^2 \left( \beta = 0 \right) - \chi^2 \left( \beta = 1 \right)$ is used calculate the rate+shape p-value, where both rate and shape information from the $\nub$ mode single-ring e-like sample and a constraint on the oscillation parameters from the other four samples (single-ring $\mu$-like and single-ring e-like, both in $\nu$ and $\nub$ modes and $\nu$ mode $\nue CC1\pi$-like) are used. Where the $\chi^2$ values are calculated by marginalising over all systematic and oscillation parameters, including the mass hierarchy.


\section{Method \& Results}

%TODO: add table of priors?
$\nueb$ appearance p-values are calculated by considering the rate-only and rate+shape test statistics of an ensemble of pseudo-experiments, generated by randomising oscillation and systematic parameters according to the priors in Tab \ref{tab:osc_param_priors_nuebarapp} \hl{[Add prior table?]} and applying statistical fluctuations. The oscillation and systematic model parameters are constrained using T2K data from four control samples ($\nu$ mode single-ring and multi-ring e-like and both $\nu$ and $\nub$ mode single-ring $\mu$-like). The pseudo-experiment ensembles are filled by sampling the oscillation and systematic parameter space around the maximum likelihood of the control samples. This ensures that the ensembles are filled with only the most likely pseudo-experiments and maintains correlations between all oscillation and systematic parameters in the four control samples.

%TODO: Do we need this paragraph?
%A right-tailed p-value is used for the $\beta = 0$ hypothesis as the $\beta = 0$ test-statistic distribution is composed solely of background events, so the right-tailed p-value gives the probability of observing a fluctuation under the $\beta = 0$ hypothesis in the background with equal to or more events than observed in the $\nub$ mode single-ring e-like sample. Similarly, a left-tailed p-value is used for the $\beta = 1$ hypothesis as the $\beta = 1$ test-statistic distribution is composed of both background and signal events, so the left-tailed p-value gives the probability of observing a fluctuation under the $\beta = 1$ hypothesis with equal to or fewer events than were observed in the $\nub$ mode single-ring e-like sample.

The distributions of the rate-only and rate+shape test statistics for each both hypotheses are shown in Fig. \ref{fig:nuebar_test_stat_dists}. These distributions are integrated from the data test statistic to obtain right(left)-tailed p-values for the $\beta = 0(1)$ hypothesis. The observed number of events in the $\nub$ mode single-ring e-like sample in Super-K was \NskObservedNueBarRunOneNineD and the observed $\Delta \chi^2$ value was \DeltachisqDataRunOneNineD. These are the rate-only and rate+shape test statistics of the data, respectively. The resulting p-values are shown in Tab. \ref{tab:data_pvals_R1-9d}. Both the rate-only and rate+shape analyses disfavour the no $\nueb$ appearance hypothesis ($\beta=0$) more than the PMNS $\nueb$ appearance hypothesis ($\beta=1$), but cannot exclude either hypothesis. The Jeffreys scale \hl{[cite]} indicates that the Bayes factor of \BayesFactorDataRunOneNineD\ corresponds to a substantial preference for the PMNS $\nueb$ appearance hypothesis over the no $\nueb$ appearance hypothesis.

\begin{figure}[!h]
	\includegraphics[width=0.5\linewidth]{figures/nuebarapp_graphs_nmargtoys=200k_nexp=20k_data_rate-only_MCsampling_VALOR_paper_Nevents.pdf}
	\hspace*{-0.5cm}\includegraphics[width=0.5\linewidth]{figures/nuebarapp_graphs_nmargtoys=200k_nexp=20k_data_rate+shape_MCsampling_VALOR_paper_Deltachisq.pdf}
	\caption{Test statistic distributions taken from the $beta=0$ and $beta=1$ pseudo-experiment ensembles for the rate-only analysis (left) and rate+shape analysis (right)}
	\label{fig:nuebar_test_stat_dists}
\end{figure}


%TODO:
%Data test-statistic
%Want RHC1Re spectrum?
%Prior table

\end{document}
